\VerdierPhysicalPageAnchor{0078}%
\VerdierSemanticAnchor{ch1-2-5-4}
\noindent\textbf{2.5.4.}\quad
The operations $\star$ and ``sum of homotopies'' have the following
properties:

\VerdierSemanticAnchor{ch1-2-5-4-eq1}
\begin{equation}
  \begin{aligned}
    u\star(b+a) &= u\star b+u\star a,\\
    (b+a)\star v &= b\star v+a\star v,\\
    uu'\star a &= u\star(u'\star a),\\
    a\star vv' &= (a\star v)\star v'.
  \end{aligned}
  \tag{2.5.4.1}
\end{equation}

\VerdierSemanticAnchor{ch1-2-5-5}
\noindent\textbf{2.5.5.}\quad
Two morphisms $f,g\colon X\rightrightarrows Y$ are said to be
\emph{homotopic} if there exists a homotopy connecting them.

\VerdierSemanticAnchor{ch1-2-5-5-difference-criterion}
Two morphisms $f,g$ are homotopic if the morphism $f-g$ is homotopic to zero.

\VerdierSemanticAnchor{ch1-2-5-5-shift-stability}
Let $f\colon X\longrightarrow Y$ be a morphism homotopic to zero. For every
$\sigma\in\mathbb Z[n]$, the morphism $T(\sigma)f$ is homotopic to zero.

\VerdierSemanticAnchor{ch1-2-5-5-null-homotopic-ideal}
The set of morphisms homotopic to zero forms an ideal of
$\operatorname{comp}^{n}(C)$, i.e. the composite, in either order, of a
morphism homotopic to zero with an arbitrary morphism is homotopic to zero,
and the direct sum of two morphisms homotopic to zero is homotopic to zero.

\VerdierSemanticAnchor{ch1-2-5-6}
\noindent\textbf{2.5.6.}\quad
Let $\varphi\colon[n]\longrightarrow[m]$ be a map.

\VerdierSemanticAnchor{ch1-2-5-6-integration-preserves-homotopy}
The functors $\int_\varphi^\Sigma$ and $\int_\varphi^\Pi$ send pairs of
homotopic morphisms of $\operatorname{comp}^{n}(C)$ to pairs of homotopic
morphisms of $\operatorname{comp}^{m}(C)$.

\VerdierSemanticAnchor{ch1-2-5-6-multiadditive-setup}
Likewise, let
$F\colon\prod_{i\in[n]}C_i\longrightarrow C$ be a multiadditive functor,

\VerdierSemanticAnchor{ch1-2-5-6-fn-functor}
\begin{equation*}
  F^{[n]}\colon
  \prod_{i\in[n]}\operatorname{comp}(C_i)
  \longrightarrow \operatorname{comp}^{n}(C)
\end{equation*}
the functor described in (2.4.2.4), let $X_i$ be objects of
$\operatorname{comp}(C_i)$ $(i\ne j)$, and let
$f,g\colon X\rightrightarrows Y$ be two homotopic morphisms of
$\operatorname{comp}(C_j)$.

\VerdierSemanticAnchor{ch1-2-5-6-fn-preserves-homotopy}
The morphisms of $\operatorname{comp}^{n}(C)$
\begin{equation*}
  F^{[n]}(X_1,\ldots,f,\ldots,X_n)
  \quad\text{and}\quad
  F^{[n]}(X_1,\ldots,g,\ldots,X_n)
\end{equation*}
are homotopic.

\VerdierSemanticAnchor{ch1-2-5-6-totalizations-preserve-homotopy}
Consequently, the morphisms of $\operatorname{comp}(C)$
\begin{equation*}
  \operatorname{comp}_{\Pi}F(X_1,\ldots,f,\ldots,X_n)
  \quad\text{and}\quad
  \operatorname{comp}_{\Pi}F(X_1,\ldots,g,\ldots,X_n)
\end{equation*}
are homotopic (the same result holds for
$\operatorname{comp}_{\Sigma}F$).

\VerdierSemanticAnchor{ch1-2-5-6-reference}
We leave it to the reader to verify all these assertions or to refer to [1].

\VerdierSemanticAnchor{ch1-2-5-7}
\noindent\textbf{2.5.7.}\quad
Let $K(C)$ denote the following category:

\VerdierSemanticAnchor{ch1-2-5-7-objects}
\noindent\hspace*{2em}--\ The objects of $K(C)$ are the objects of $\operatorname{comp}(C)$.
